\documentclass[a4paper,12pt]{article}
\usepackage[utf8]{inputenc}
\usepackage[T1]{fontenc}
\usepackage[italian]{babel}
\usepackage{listings}
\usepackage{xcolor}
\usepackage{hyperref}
\usepackage[margin=3cm]{geometry}

\title{Documentazione Progetto Sistemi Operativi: PandOSsh Fase 1}
\author{Luca Argentino, Milo Disalvatore, Matteo Mazzetti, Zeno Maccari}
\begin{document}


\maketitle

\section{Gestione PCB (pcb.c)}
Il modulo \texttt{pcb.c} si occupa della gestione dei \textbf{Process Control Blocks (PCB)}, occupandosi della loro allocazione/deallocazione e del mantenimento delle strutture a coda e ad albero.

\vspace{0.5cm}
\subsection{Strutture Dati Globali}
Il modulo utilizza le seguenti variabili statiche globali per gestire il pool di PCB:

\begin{itemize}
    \item \texttt{static pcb\_t pcbFree\_table[MAXPROC]}: Array statico che contiene i PCB disponibili per il sistema.
    \item \texttt{static struct list\_head pcbFree\_h}: Testa della lista che collega i PCB attualmente non in uso (liberi).
    \item \texttt{static int next\_pid}: Contatore intero utilizzato per assegnare un Identificatore di Processo (PID) unico a ogni nuovo PCB allocato.
\end{itemize}

\vspace{0.5cm}
\subsection{Descrizione delle Funzioni di Allocazione}

\subsubsection{\texttt{void initPcbs()}}
\begin{enumerate}
    \item Inizializza la testa della lista dei PCB liberi (\texttt{pcbFree\_h}) tramite la macro \texttt{INIT\_LIST\_HEAD}.
    \item Itera attraverso l'array statico \texttt{pcbFree\_table} e inserisce ogni elemento nella lista \texttt{pcbFree\_h} utilizzando il campo \texttt{p\_list}.
\end{enumerate}

\subsubsection{\texttt{void freePcb(pcb\_t *p)}}
\begin{enumerate}
    \item Inserisce il PCB puntato da \texttt{p} nella lista dei PCB liberi (\texttt{pcbFree\_h}), rendendolo disponibile per future allocazioni.
\end{enumerate}

\subsubsection{\texttt{pcb\_t *allocPcb()}}
\begin{enumerate}
    \item Controlla se la lista \texttt{pcbFree\_h} è vuota. Se sì, ritorna \texttt{NULL}.
    \item Se è presente un PCB libero:
    \begin{itemize}
        \item Rimuove il primo elemento dalla lista \texttt{pcbFree\_h}.
        \item Inizializza tutti i campi del PCB a \texttt{NULL} o \texttt{0} (inclusi i puntatori di lista e di stato).
        \item Assegna al campo \texttt{p\_pid} il valore corrente di \texttt{next\_pid} e incrementa quest'ultimo.
        \item Ritorna il puntatore al PCB inizializzato.
    \end{itemize}
\end{enumerate}

\vspace{0.5cm}
\subsection{Descrizione delle Funzioni per Code di Processi} %-----PARTE DI MAZZO
\subsubsection{\texttt{int emptyProcQ(struct list\_head *head)}}
\begin{enumerate}
    \item Ritorna \texttt{TRUE} se la lista puntata da \texttt{head} è vuota, \texttt{FALSE} altrimenti.
\end{enumerate}

\subsubsection{\texttt{void insertProcQ(struct list\_head *head, pcb\_t *p)}}
\begin{enumerate}
    \item Memorizza la priorità del PCB puntato da \texttt{p}.
    \item Scorre tutta la lista associata ad \texttt{head} e confronta le priorità dei PCB che incontra con quella memorizzata.
    \item Appena incontra un PCB con priorità minore di quella di \texttt{p} inserisce il PCB \texttt{p}.
\end{enumerate}
\subsubsection{\texttt{pcb\_t *headProcQ(struct list\_head *head)}}
\begin{enumerate}
    \item Controlla se l'elemento sentinella \texttt{head} punta ad una coda nulla e in tal caso ritorna \texttt{NULL}.
    \item Se così non fosse allora restituisce il primo PCB della lista.
\end{enumerate}

\subsubsection{\texttt{pcb\_t *removeProcQ(struct list\_head *head)}}
\begin{enumerate}
    \item Se \texttt{head} punta ad una coda vuota allora viene ritornato \texttt{NULL}.
    \item Altrimenti salva temporaneamente il primo PCB per ritornarlo.
    \item Rimuove il primo PCB.
    \item Restituisce il PCB appena salvato, ossia quello rimosso. 
\end{enumerate}

\subsubsection{\texttt{pcb\_t *outProcQ(struct list\_head *head, pcb\_t *p)}}
\begin{enumerate}
    \item Scorre tutta la lista puntata da \texttt{head} alla ricerca del PCB di cui \texttt{p} è il puntatore.
    \begin{itemize}
        \item Se trova p allora lo rimuove e lo restituisce.
        \item Se p non è presente nella lista allora ritorna \texttt{NULL}.
    \end{itemize}
\end{enumerate}


\vspace{0.5cm}
\subsection{Descrizione delle Funzioni per Alberi di Processi}

\subsubsection{\texttt{int emptyChild(pcb\_t *p)}}
\begin{enumerate}
    \item Verifica se il PCB puntato da \texttt{p} ha dei figli controllando se la lista \texttt{p\_child} è vuota.
    \item Ritorna \texttt{TRUE} se la lista è vuota (o se \texttt{p} è \texttt{NULL}), altrimenti ritorna \texttt{FALSE}.
\end{enumerate}

\subsubsection{\texttt{void insertChild(pcb\_t *prnt, pcb\_t *p)}}
\begin{enumerate}
    \item Se entrambi i puntatori sono validi, imposta \texttt{prnt} come padre del PCB \texttt{p} (campo \texttt{p\_parent}).
    \item Inserisce il PCB \texttt{p} nella lista dei figli del padre (\texttt{prnt->p\_child}) utilizzando il campo \texttt{p\_sib}.
\end{enumerate}

\subsubsection{\texttt{pcb\_t *removeChild(pcb\_t *p)}}
\begin{enumerate}
    \item Controlla se il PCB \texttt{p} ha figli tramite \texttt{emptyChild}. Se non ne ha, ritorna \texttt{NULL}.
    \item Altrimenti, individua il primo figlio nella lista \texttt{p\_child}.
    \item Rimuove tale figlio dalla lista e ne imposta il puntatore al padre (\texttt{p\_parent}) a \texttt{NULL}.
    \item Ritorna il puntatore al PCB del figlio rimosso.
\end{enumerate}

\subsubsection{\texttt{pcb\_t *outChild(pcb\_t *p)}}
\begin{enumerate}
    \item Verifica se il PCB \texttt{p} ha effettivamente un padre. Se non lo ha, ritorna \texttt{NULL}.
    \item Rimuove il PCB \texttt{p} dalla lista dei fratelli (e quindi dalla lista dei figli del padre) tramite il campo \texttt{p\_sib}.
    \item Resetta il puntatore \texttt{p\_parent} a \texttt{NULL}.
    \item Ritorna il puntatore \texttt{p}.
\end{enumerate}


\section{Gestione Semafori (asl.c)}
Il modulo \texttt{asl.c} implementa la gestione della \textbf{Active Semaphore List (ASL)}.

La struttura dati ASL è utilizzata per la gestione della lista dei semafori attivi, ovvero quei semafori in cui almeno un PCB è bloccato.

\vspace{0.5cm}
\subsection{Strutture Dati Globali}
Il modulo utilizza le seguenti variabili statiche globali per gestire lo stato dei semafori:

\begin{itemize}
    \item \texttt{static semd\_t semd\_table[MAXPROC]}: Array statico contente i semafori liberi con dimensione massima \texttt{(MAXPROC)}.
    \item \texttt{static struct list\_head semdFree\_h}: Testa della lista che collega i descrittori di semafori liberi.
    \item \texttt{static struct list\_head semd\_h}: Testa della lista che rappresenta la ASL che contiene i descrittori dei semafori attivi.
\end{itemize}

\vspace{0.5cm}
\subsection{Descrizione delle Funzioni}

\subsubsection{\texttt{void initASL()}}
\begin{enumerate}
    \item Inizializza le teste delle liste \texttt{semdFree\_h} e \texttt{semd\_h} utilizzando la macro \texttt{INIT\_LIST\_HEAD}.
    \item Itera attraverso l'array \texttt{semd\_table} e inserisce ogni elemento nella lista dei descrittori liberi (\texttt{semdFree\_h}) utilizzando il campo \texttt{s\_link}.
\end{enumerate}

\subsubsection{\texttt{int insertBlocked(int* semAdd, pcb\_t* p)}}


\begin{enumerate}
    \item Scorre la ASL (\texttt{semd\_h}) per cercare un semaforo con chiave uguale a \texttt{semAdd}.
    \item \textbf{Se il semaforo esiste:}
    \begin{itemize}
        \item Inserisce il PCB \texttt{p} in coda alla lista processi del semaforo (\texttt{s\_procq}).
        \item Aggiorna il puntatore \texttt{p->p\_semAdd} del PCB.
        \item Ritorna \texttt{FALSE}.
    \end{itemize}
    \item \textbf{Se il semaforo non esiste:}

    Controlla se la lista dei liberi \texttt{semdFree\_h} è vuota. Se sì, ritorna \texttt{TRUE}.

    Altrimenti: 
    \begin{itemize}
        
        \item Rimuove il primo semaforo  disponibile dalla \texttt{semdFree\_h}.
        \item Inizializza il nuovo semaforo: imposta la chiave \texttt{s\_key}, inizializza la coda processi \texttt{s\_procq}.
        \item Aggiunge il nuovo semaforo in coda alla ASL (\texttt{semd\_h}).
        \item Aggiunge il PCB \texttt{p} alla coda del nuovo semaforo e aggiorna il riferimento \texttt{p->p\_semAdd}.
        \item Ritorna \texttt{FALSE}.
    \end{itemize}
\end{enumerate}

\subsubsection{\texttt{pcb\_t* removeBlocked(int* semAdd)}}


\begin{enumerate}
    \item Cerca il semaforo con chiave \texttt{semAdd} nella ASL.
    \item Se non trovato, ritorna \texttt{NULL}.
    \item Se trovato:
    \begin{itemize}
        \item Identifica e rimuove il primo PCB nella coda \texttt{s\_procq}.
        \item Controlla se la coda processi del semaforo è diventata vuota (\texttt{emptyProcQ}).
        \begin{itemize}
            \item Se la coda è vuota, il semaforo viene disattivato: viene rimosso dalla ASL e reinserito nella lista dei liberi 
        \texttt{semdFree\_h}.
        \end{itemize}
        \item Ritorna il puntatore al PCB rimosso.
    \end{itemize}
\end{enumerate}

\subsubsection{\texttt{pcb\_t* outBlocked(pcb\_t* p)}}


\begin{enumerate}
    \item Scorre la ASL per trovare il semaforo corrispondente a \texttt{p->p\_semAdd}.
    \item Una volta trovato il semaforo, scorre la sua coda processi (\texttt{s\_procq}) per trovare il PCB \texttt{p}.
    \item Se il PCB viene trovato:
    \begin{itemize}
        \item Rimuove il PCB dalla lista.
        \item Verifica se la coda del semaforo è diventata vuota. Se sì, rimuove il descrittore del semaforo dalla ASL e lo restituisce alla lista dei liberi.
        \item Ritorna il puntatore a \texttt{p}.
    \end{itemize}
    \item Se il semaforo non viene trovato o il PCB non è nella coda, ritorna \texttt{NULL}.
\end{enumerate}

\subsubsection{\texttt{pcb\_t* headBlocked(int* semAdd)}}


\begin{enumerate}
    \item Cerca il semaforo con chiave \texttt{semAdd} nella ASL.
    \item Se trovato, utilizza la funzione \texttt{headProcQ} (definita in \texttt{pcb.c}) per restituire il primo elemento della coda processi.
    \item Se il semaforo non esiste, ritorna \texttt{NULL}.
\end{enumerate}

\end{document}
